\section{Aprendizaje por refuerzo y offline RL: de dónde vienen los datos}

Esta sección presenta los antecedentes de RL mínimos, pero completos. El lector necesita distinguir entre state, action, reward, return, policy y dataset support. Si estos conceptos se confunden, en lo que sigue es fácil tomar el return-to-go por un futuro conocido o suponer que el offline RL elimina por completo el problema de la exploración.

\subsection{MDP, policy y cumulative return}

Un Markov decision process (MDP) estándar puede escribirse como $\mathcal{M}=(\mathcal{S},\mathcal{A},P,r,\gamma)$. En el estado $s_t\in\mathcal{S}$, la policy $\pi(a_t\mid s_t)$ elige la acción $a_t\in\mathcal{A}$; el entorno transita según el transition kernel $P(s_{t+1}\mid s_t,a_t)$ y produce la reward $r_t$. El objetivo del aprendizaje es maximizar el retorno descontado esperado:
\[
J(\pi)=\mathbb{E}_{\tau\sim\pi}\left[\sum_{t=1}^{T}\gamma^{t-1}r_t\right],
\qquad 0<\gamma\le 1.
\]
Aquí $\tau=(s_1,a_1,r_1,\ldots,s_T,a_T,r_T)$ es la trajectory; $\gamma$ es el discount factor, que sirve para reducir el peso de las recompensas lejanas o para garantizar que la suma en horizonte infinito sea estable.

\lecturefigure{slide-05-background-reinforcement-learning.jpg}{Ciclo cerrado de RL: el agent elige una action según el state y el environment devuelve la reward y el siguiente estado.}{Video oficial de Stanford Online, 04:00.}

\begin{knowledgebox}{Términos clave: las cantidades que más se confunden en RL}
\begin{tabularx}{\textwidth}{L{2.8cm} X X}
\toprule
Término & Problema que resuelve & Relación con esta clase \\
\midrule
reward $r_t$ & Retroalimentación escalar obtenida en el paso actual & Después del rollout se usa para actualizar el return-to-go restante \\
return $G_t$ & Reward futura acumulada a partir del instante $t$ & En los datos de entrenamiento puede recalcularse a partir de la trayectoria completa \\
policy $\pi$ & Cómo elegir una action dada la información disponible & Decision Transformer implementa la policy de forma implícita mediante action prediction \\
value / critic & Estima el valor a largo plazo de un estado o de una acción & El método original no depende de un critic para entrenar al actor, pero más adelante se muestra que puede cambiarse a un critic target \\
credit assignment & Determinar cuánta responsabilidad tienen las acciones tempranas en los resultados lejanos & Key-to-Door lo pone a prueba con una reward dispersa entre etapas \\
\bottomrule
\end{tabularx}
\end{knowledgebox}

\subsection{El conjunto de datos fijo del offline RL}

El RL en línea puede interactuar con el entorno de forma continua durante el entrenamiento; el offline RL (aprendizaje por refuerzo fuera de línea) solo usa un conjunto de datos fijo recolectado de antemano
\[
\mathcal{D}=\{\tau^{(i)}\}_{i=1}^{N},
\]
cuyas trayectorias pueden provenir de una policy aleatoria, de una policy entrenada a medias, de una policy experta o de una mezcla de varias policies. Los datos fijos reducen el costo y el riesgo de la interacción en robots reales, en sistemas médicos o industriales, pero plantean un problema más agudo: durante el rollout, el modelo puede visitar regiones state-action que el conjunto de datos casi no cubre.

\lecturefigure{slide-06-background-offline-rl.jpg}{El offline RL solo usa logged interactions registradas de antemano y no realiza nuevos trial-and-error.}{Video oficial de Stanford Online, 05:20.}

\begin{importantbox}{Dataset support (soporte del conjunto de datos)}
Si el conjunto de datos casi no tiene muestras de la acción $a$ cerca de cierto estado $s$, es difícil estimar con confiabilidad el resultado de ejecutar esa acción. El modelo puede producir una action sintácticamente razonable, pero la aproximación de funciones no crea por sí sola evidencia contrafactual. Lo esencial del offline RL no es que baste con «tener un conjunto de datos grande», sino que los datos cubran la distribución de comportamientos y estados que se necesita en el despliegue.
\end{importantbox}

\begin{warningbox}{Los dos niveles de la distribution shift}
El primer nivel es que los datos de entrenamiento los genera la behavior policy $\mu$, mientras que en el despliegue se usa la learned policy $\pi$; el segundo nivel es que un pequeño error en una acción empuja al agent hacia un estado nuevo, y las predicciones siguientes operan sobre entradas cada vez menos conocidas. Decision Transformer evita el Bellman backup explícito, pero no hace que estos dos niveles de desplazamiento de distribución desaparezcan por sí solos.
\end{warningbox}

\subsection{Por qué el expositor dice que el RL «no escala»}

Hacia 2021, los grandes language models ya alcanzaban de decenas a cientos de miles de millones de parámetros, mientras que las RL policies habituales solían ser mucho más pequeñas y su entrenamiento mostraba fluctuaciones más marcadas. La diapositiva atribuye la diferencia a la inestabilidad del RL objective; este juicio debe entenderse como la motivación de la clase: el bootstrap target, el on-policy distribution change, la exploration y la reward scale aumentan la dificultad de la optimización. No quiere decir que el número de parámetros por sí mismo mida la inteligencia de una policy.

\lecturefigure{slide-07-motivating-challenge-rl-does-not-scale.jpg}{Motivación de la clase: los grandes modelos de lenguaje de la época eran mucho mayores que las RL policies habituales, y su entrenamiento era más estable.}{Video oficial de Stanford Online, 07:00.}

\begin{knowledgebox}{Lectura de la figura: primero comparar la interfaz de entrenamiento, después el número de parámetros}
El texto de la izquierda compara la model scale y la curva de la derecha compara la training stability. El argumento que realmente respalda a Decision Transformer es este: si la action prediction puede usar MSE o cross-entropy estándar, entonces pueden reutilizarse las herramientas de optimización y regularización del aprendizaje supervisado. La figura no demuestra que «cuanto mayor es el modelo, mejor es el control», ni controla el entorno, la cantidad de datos, el action space ni el presupuesto de evaluación; por lo tanto, la diferencia en parámetros no puede interpretarse como una relación causal única.
\end{knowledgebox}

\teachervoice{El docente enfatiza que, en los offline benchmarks, los datos suelen provenir del replay buffer de un agent de RL en línea, por ejemplo, las interacciones guardadas cuando el entrenamiento alcanzó un nivel intermedio. La interfaz del algoritmo puede ser agnostic respecto al origen de los datos, pero las conclusiones experimentales no pueden ignorar la collection policy: la calidad y la cobertura de los datos influyen directamente en lo que el modelo puede aprender.}

\subsection{Resumen de la sección}

El objetivo del RL es maximizar el retorno acumulado en el entorno real; el offline RL, en cambio, aprende una policy a partir de logged trajectories fijas. Los datos fijos eliminan las interacciones nuevas durante el entrenamiento, pero agravan los problemas de dataset support y de distribution shift. Lo que Decision Transformer cambia a continuación es la formulación del entrenamiento, no estas restricciones estadísticas.
